\documentclass[journal]{IEEEtran}

% ── Packages ──────────────────────────────────────────────────────
\usepackage{cite}
\usepackage{amsmath,amssymb,amsfonts}
\usepackage{algorithmic}
\usepackage{graphicx}
\usepackage{textcomp}
\usepackage{xcolor}
\usepackage{booktabs}
\usepackage{hyperref}
\usepackage{multirow}
\usepackage{array}
\usepackage{url}
\usepackage{microtype}

% ── Title ─────────────────────────────────────────────────────────
\title{Aletheia: An Offline-First Clinical Decision Support
System for Differential Diagnosis in Low-Resource Healthcare Settings}

\author{Joseph Walusimbi,~\IEEEmembership{}
  Ann Move Oguti,~\IEEEmembership{}
  Abubakhari Sserwadda,~\IEEEmembership{}
  Precious Boss Kasasira,~\IEEEmembership{}
  and~Charles Brian Okoboi~\IEEEmembership{}

\thanks{J. Walusimbi, A. M. Oguti, and A. Sserwadda are with the
Department of Electronics and Computer Engineering, Soroti University,
Soroti, Uganda (e-mail: 2401600068@sun.ac.ug; amoguti@sun.ac.ug;
asserwadda@sun.ac.ug).}

\thanks{P. B. Kasasira and C. B. Okoboi are with the School of Health
Sciences, Soroti University, Soroti, Uganda
(e-mail: kasasirakp40@gmail.com; charlesbryahn@gmail.com).}}

\begin{document}

\maketitle

% ══════════════════════════════════════════════════════════════════
\begin{abstract}
Access to specialist clinical expertise remains severely limited
across sub-Saharan Africa, where physician-to-patient ratios can
fall below 1:25{,}000 in rural settings. Existing AI-assisted
diagnostic tools predominantly require reliable internet connectivity
and high-specification hardware, rendering them impractical for
frontline healthcare workers in district hospitals and health
centres. This paper presents \textit{Aletheia}, an offline-first
clinical decision support system that generates ranked differential
diagnoses for low-resource healthcare contexts across sub-Saharan
Africa. Aletheia is built upon Qwen2.5-3B-Instruct, fine-tuned using
Quantised Low-Rank Adaptation (QLoRA) on a curated dataset of
27{,}000 clinical reasoning samples spanning 50 disease conditions
with elevated prevalence in East Africa. The system runs entirely
on commodity laptop hardware (Intel Core i5, 8\,GB DDR4 RAM) without
internet connectivity, occupying only 1.80\,GB of storage in its
primary deployment format. Evaluation demonstrates a Top-1
differential diagnosis accuracy of 80.0\%, Top-3 accuracy of
100.0\%, BERTScore-F1 of 0.909, and METEOR of 0.467 across ten
representative clinical case categories. The system achieves an
Expected Calibration Error (ECE) of 0.275 and passes the Africa
Deep Tech Challenge 2026 (ADTC 2026) memory budget constraint of
7\,168\,MB, achieving a peak inference RAM of approximately
3\,630\,MB on the standardised benchmark laptop. These results
demonstrate the feasibility of deploying large language model-based
differential diagnosis at the primary care level in
resource-constrained settings without cloud infrastructure.
\end{abstract}

\begin{IEEEkeywords}
Clinical decision support, large language models, QLoRA fine-tuning,
offline inference, sub-Saharan Africa, low-resource healthcare,
differential diagnosis, GGUF quantisation, East Africa.
\end{IEEEkeywords}

% ══════════════════════════════════════════════════════════════════
\section{Introduction}
\label{sec:intro}

\IEEEPARstart{T}{he} burden of disease in sub-Saharan Africa is
disproportionately high relative to the availability of clinical
expertise. Uganda, with a population exceeding 48 million, has a
physician density of approximately 0.17 per 1{,}000 population
\cite{who2023workforce}---among the lowest globally. In rural and
peri-urban settings such as Soroti District in Eastern Uganda, a
single clinical officer may be responsible for 80--120 patient
consultations per day, leaving fewer than five minutes per patient
for history-taking, examination, differential diagnosis, and
investigation planning. Under these conditions, diagnostic errors
and missed critical presentations are an inevitable consequence of
cognitive overload rather than clinical incompetence.

Artificial intelligence-assisted clinical decision support systems
(CDSS) have demonstrated considerable promise in high-income
settings, with models achieving diagnostic accuracy comparable to
specialist physicians in radiology \cite{rajpurkar2022ai},
dermatology \cite{esteva2017dermatologist}, and general internal
medicine \cite{mcdermott2021comprehensive}. However, the vast
majority of such systems require persistent internet connectivity,
cloud-based inference infrastructure, and hardware specifications
far beyond what is available in sub-Saharan African health
facilities \cite{nkosi2023challenges}.

Recent advances in large language model (LLM) compression,
particularly 4-bit quantisation via the GGUF format and low-rank
adaptation (LoRA) fine-tuning, have created a new possibility:
deploying clinically useful AI reasoning on commodity hardware
without internet connectivity \cite{dettmers2023qlora,gerganov2023llamacpp}.
The combination of these techniques enables a 3-billion parameter
language model to be compressed to under 2\,GB and executed on a
standard laptop at practical inference speeds.

This paper presents \textit{Aletheia} (from the Greek for
\textit{truth} or \textit{disclosure}), an offline-first clinical
decision support system that:
\begin{itemize}
  \item Runs entirely on-device with no internet dependency,
        suitable for deployment at district hospital and health
        centre level across sub-Saharan Africa;
  \item Provides ranked differential diagnoses with probability
        estimates, evidence-based investigation recommendations,
        clinical rationale, red flag identification, and follow-up
        question generation;
  \item Is fine-tuned on a dataset weighted toward African disease
        epidemiology --- with particular emphasis on East African
        prevalence patterns --- including conditions rarely represented
        in publicly available medical AI benchmarks;
  \item Satisfies the memory constraints of the Africa Deep Tech
        Challenge 2026 (ADTC 2026) \cite{adtc2026}, which requires
        models to operate within 7\,168\,MB on a standardised laptop
        (Intel Core i5 10th--12th generation, 8\,GB DDR4,
        Ubuntu\,22.04, no discrete GPU), with Aletheia achieving a
        peak RAM requirement of approximately 3\,730\,MB --- a margin
        of 3\,438\,MB below the ceiling.
\end{itemize}

The remainder of this paper is structured as follows.
Section~\ref{sec:related} reviews related work in medical LLMs
and clinical decision support for low-resource settings.
Section~\ref{sec:dataset} describes the training dataset.
Section~\ref{sec:method} presents the model architecture, fine-tuning
methodology, and deployment pipeline.
Section~\ref{sec:experiments} details the experimental setup and
evaluation metrics.
Section~\ref{sec:results} reports quantitative results.
Section~\ref{sec:discussion} discusses findings, limitations, and
future directions.
Section~\ref{sec:conclusion} concludes the paper.

% ══════════════════════════════════════════════════════════════════
\section{Related Work}
\label{sec:related}

\subsection{Large Language Models in Clinical Medicine}

The application of transformer-based language models to clinical
reasoning has accelerated substantially since the release of
GPT-4 \cite{openai2023gpt4} and its successors. Med-PaLM\,2
\cite{singhal2023medpalm2} achieved expert-level performance on
USMLE-style medical questions, while BioMedLM \cite{bolton2022biomedlm}
and ClinicalBERT \cite{huang2019clinicalbert} demonstrated the value
of domain-specific pre-training. However, these systems are
predominantly evaluated on US or European clinical benchmarks and
require cloud-based inference infrastructure, limiting their
applicability in resource-constrained settings.

\subsection{Clinical Decision Support in Low-Resource Settings}

Prior work on CDSS for sub-Saharan Africa has focused primarily on
rule-based expert systems \cite{fraser2007clinician}, mobile health
(mHealth) applications \cite{labrique2013mhealth}, and telemedicine
platforms \cite{bastawrous2015ehealth}. While these approaches have
demonstrated utility, they either lack the reasoning depth of modern
LLMs or depend on internet connectivity for core functionality.
Symptom-checker applications such as Ada \cite{ada2020} and Babylon
\cite{babylon2020} provide structured differential diagnosis but
require continuous data transmission to cloud inference endpoints.

\subsection{Efficient LLM Deployment}

The quantisation of large language models for edge deployment has
been advanced by work on GPTQ \cite{frantar2022gptq},
AWQ \cite{lin2023awq}, and the GGUF format supported by
llama.cpp \cite{gerganov2023llamacpp}. QLoRA \cite{dettmers2023qlora}
demonstrated that high-quality fine-tuning of quantised models is
achievable with minimal hardware, enabling fine-tuning of 7B+
parameter models on consumer GPUs. The combination of QLoRA
fine-tuning and GGUF deployment enables a workflow where training
occurs on available cloud hardware and the resulting model is
deployed on commodity local hardware.

\subsection{African Medical AI}

Dedicated African medical AI datasets and benchmarks remain scarce.
WHO AFRO health statistics \cite{afrihealth2023} and regional
epidemiological reports provide limited structured clinical AI benchmark data. The MedQA \cite{jin2021medqa} and
MedMCQA \cite{pmlr-v174-pal22a} datasets offer large-scale medical
question-answering data but are predominantly US and Indian in
clinical context respectively. Aletheia addresses this gap through
a purpose-built synthetic dataset weighted to reflect East African
disease epidemiology.

% ══════════════════════════════════════════════════════════════════
\section{Dataset}
\label{sec:dataset}

\subsection{Dataset Composition}

The Aletheia training dataset was constructed from three sources:

\begin{enumerate}
  \item \textbf{Aletheia-Synthetic}: 20{,}000 structured clinical
        reasoning samples generated from 50 hand-crafted clinical
        case templates covering conditions prevalent in East Africa.
  \item \textbf{MedQA-USMLE} \cite{jin2021medqa}: 9{,}358 filtered
        questions from the MedQA dataset, selected for relevance to
        tropical medicine, infectious disease, obstetrics, and
        paediatrics using a 180-term keyword filter.
  \item \textbf{MedMCQA} \cite{pmlr-v174-pal22a}: 10{,}000 filtered
        questions from the MedMCQA dataset, selected using the same
        keyword filter and capped at 10{,}000 samples.
\end{enumerate}

After applying a 60/20/20 mixing ratio across all three sources
and a minimum-token quality filter (50 tokens), the final dataset
comprised 27{,}000 training samples and 3{,}000 evaluation samples,
totalling 30{,}000 samples --- an increase of 25\% over the
initial single-source configuration.
Table~\ref{tab:dataset} summarises the dataset statistics.

\begin{table}[h!]
\centering
\caption{Training Dataset Statistics}
\label{tab:dataset}
\begin{tabular}{lc}
\toprule
\textbf{Property} & \textbf{Value} \\
\midrule
Total samples            & 30{,}000 \\
Training samples         & 27{,}000 \\
Evaluation samples       & 3{,}000 \\
Synthetic samples        & 18{,}000 (60.0\%) \\
MedQA-USMLE samples      & 6{,}000 (20.0\%) \\
MedMCQA samples          & 6{,}000 (20.0\%) \\
Clinical conditions      & 50 \\
Reasoning task types     & 8 \\
Total tokens             & $\sim$4{,}000{,}000 \\
Mean tokens per sample   & 133 \\
\bottomrule
\end{tabular}
\end{table}

\subsection{Clinical Conditions}

The 50 clinical conditions were selected to reflect the East African
disease burden across eight categories:
infectious and tropical disease (12 conditions),
respiratory (3), cardiovascular (3), obstetric and gynaecological
(4), paediatric (4), neurological (2), renal and endocrine (4),
surgical and trauma (3), and other specialties (15).
Conditions were weighted by estimated incidence in East Africa
(africa\_weight parameter, range 3--5), with malaria, HIV/AIDS,
tuberculosis, cerebral malaria, eclampsia, severe acute malnutrition,
and neonatal sepsis receiving the highest weights.

\subsection{Reasoning Task Types}

Each clinical case was instantiated across eight reasoning task
types to teach multi-step clinical reasoning rather than simple
diagnosis recall:

\begin{enumerate}
  \item \textit{Initial differential diagnosis} --- ranked
        differentials from presenting symptoms;
  \item \textit{Test recommendation} --- evidence-based
        investigation prioritisation;
  \item \textit{Evidence update} --- Bayesian probability revision
        after test results;
  \item \textit{Rationale explanation} --- plain-language clinical
        reasoning;
  \item \textit{Follow-up questions} --- most discriminating next
        question generation;
  \item \textit{Severity assessment} --- acuity classification and
        level-of-care determination;
  \item \textit{Treatment hint} --- district-hospital-level
        immediate management;
  \item \textit{Red flag identification} --- immediate
        escalation triggers.
\end{enumerate}

% ══════════════════════════════════════════════════════════════════
\section{Methodology}
\label{sec:method}

\subsection{Base Model Selection}

Qwen2.5-3B-Instruct \cite{qwen2025} was selected as the base model
for the following reasons:
\begin{itemize}
  \item Parameter count (3.09B) is sufficient for multi-step clinical
        reasoning while remaining within the ADTC memory budget after
        quantisation;
  \item Native bfloat16 support enables stable training on A100
        hardware without numerical instability;
  \item Strong baseline performance on instruction-following tasks
        reduces the fine-tuning sample requirement;
  \item Permissive licensing supports open deployment in public
        health contexts.
\end{itemize}

\subsection{Fine-Tuning with QLoRA}

We applied Quantised Low-Rank Adaptation (QLoRA) \cite{dettmers2023qlora}
to fine-tune Qwen2.5-3B-Instruct. On the A100 training hardware,
we used full bfloat16 precision without quantisation to maximise
training quality. The LoRA adapters were attached to all linear
projection layers: \texttt{q\_proj}, \texttt{k\_proj},
\texttt{v\_proj}, \texttt{o\_proj}, \texttt{gate\_proj},
\texttt{up\_proj}, and \texttt{down\_proj}.

Table~\ref{tab:config} summarises the training configuration.

\begin{table}[h!]
\centering
\caption{Experimental Configuration}
\label{tab:config}
\begin{tabular}{ll}
\toprule
\textbf{Parameter} & \textbf{Value} \\
\midrule
Base model              & Qwen2.5-3B-Instruct \\
Total parameters        & 3.09B \\
Trainable (LoRA)        & 59{,}867{,}136 (1.94\%) \\
LoRA rank ($r$)         & 32 \\
LoRA alpha ($\alpha$)   & 64 \\
LoRA dropout            & 0.05 \\
Target modules          & All 7 linear projections \\
Training epochs         & 3 \\
Batch size (per device) & 8 \\
Gradient accumulation   & 2 \\
Effective batch size    & 16 \\
Learning rate           & $2 \times 10^{-4}$ \\
LR scheduler            & Cosine with warmup \\
Warmup ratio            & 0.05 \\
Weight decay            & 0.01 \\
Sequence length         & 1{,}024 tokens \\
Optimiser               & AdamW \\
Precision               & BFloat16 \\
Training hardware       & NVIDIA A100-SXM4-80GB \\
Training time           & 1.92 hours \\
\bottomrule
\end{tabular}
\end{table}

\subsection{Deployment Pipeline}

After fine-tuning, the LoRA adapters were merged into the base
model weights using the PEFT library \cite{peft2023}. The merged
model was then converted to the GGUF format using llama.cpp
\cite{gerganov2023llamacpp} via a two-step process: conversion to
16-bit floating point GGUF (6.18\,GB), followed by quantisation
to Q4\_K\_M (1.80\,GB) and Q2\_K (1.19\,GB) using
\texttt{llama-quantize}. The primary deployment target is the
Q4\_K\_M quantisation, which achieves approximately 98\% of F16
quality at 31\% of the file size.

The system provides two interfaces for clinical use: a terminal CLI for scripting and integration, and a web-based graphical user interface (GUI) built with Gradio \cite{aletheia_github} that runs locally in the browser without internet connectivity. Both interfaces accept structured clinical input and return structured JSON containing ranked differentials, recommended tests, rationale, and red flags.

\subsection{Offline-First Architecture}

Aletheia implements a local-first, sync-when-available architecture:
\begin{itemize}
  \item \textbf{Offline mode}: Full diagnostic capability, no
        degradation. The 1.80\,GB GGUF model and inference engine
        operate entirely on-device;
  \item \textbf{Connected mode}: Optional model update download,
        anonymised aggregate usage logging for quality improvement;
  \item \textbf{Update safety}: New model versions are validated on
        a held-out test set before activation; rollback to the
        previous version is retained for 30 days.
\end{itemize}

% ══════════════════════════════════════════════════════════════════
\section{Experimental Setup}
\label{sec:experiments}

\subsection{Training Infrastructure}

Training was performed on Google Colab Pro using an NVIDIA
A100-SXM4-80GB GPU (85.1\,GB VRAM) with PyTorch 2.11.0,
Transformers 4.44.2, PEFT 0.12.0, TRL 0.10.1, and
Accelerate 0.34.2. The complete training run required 1.92 hours
for 4{,}050 steps across 3 epochs.

\subsection{Evaluation Protocol}

Model performance was evaluated across four dimensions:

\subsubsection{Clinical Accuracy}
Top-$k$ accuracy was computed by checking whether the correct
diagnosis appeared in the model's top-1 or top-3 ranked
differentials. Per-condition precision, recall, and F1 were
computed over the 10 core evaluation case categories. Accuracy
was further stratified by clinical severity (Critical, High,
Moderate, Low) and reasoning task type.

\subsubsection{Language Quality}
Text generation quality was assessed using ROUGE-1, ROUGE-2,
and ROUGE-L \cite{lin2004rouge} (F1 scores), BERTScore-F1
\cite{zhang2019bertscore} using the \texttt{roberta-large} encoder,
and METEOR \cite{banerjee2005meteor}.

\subsubsection{Calibration}
Probability calibration was assessed using Expected Calibration
Error (ECE) \cite{naeini2015ece}, Maximum Calibration Error (MCE),
and Brier score \cite{brier1950brier} stratified by clinical
severity. A reliability diagram was constructed using 10 equal-width
bins.

\subsubsection{Baseline Comparison}
The fine-tuned Aletheia model was compared against the unmodified
Qwen2.5-3B-Instruct base model (zero-shot) on Top-1 accuracy and
ROUGE-1 to quantify the contribution of fine-tuning.

\subsection{Hardware Compliance Testing}

Inference memory requirements were estimated against the ADTC 2026
standard laptop specification (Intel Core i5 10th--12th generation,
8\,GB DDR4, Ubuntu\,22.04, no discrete GPU) using the formula:

\begin{equation}
  M_{\text{total}} = M_{\text{OS}} + M_{\text{model}} + M_{\text{KV}} + M_{\text{runtime}} + M_{\text{app}}
  \label{eq:memory}
\end{equation}

where $M_{\text{OS}} = 900$\,MB, $M_{\text{KV}} = 400$\,MB
(1{,}024 token context), $M_{\text{runtime}} = 300$\,MB, and
$M_{\text{app}} = 200$\,MB.

% ══════════════════════════════════════════════════════════════════
\section{Results}
\label{sec:results}

\subsection{Training Dynamics}

Fig.~\ref{fig:training} shows the training and validation loss
curves, perplexity, learning rate schedule, and gradient norm over
4{,}050 training steps. Training loss decreased from an initial
value of approximately 1.33 in early steps to a final value of
0.5197. The relatively higher terminal loss compared to
single-source training runs reflects the increased task diversity
introduced by the MedMCQA component, which exposes the model to
a broader range of clinical question styles and reasoning patterns.
Crucially, this broader training signal produced substantial gains
in diagnostic accuracy, with Top-1 accuracy improving from 70.0\%
to 80.0\% and Top-3 accuracy reaching 100.0\% --- demonstrating
that loss alone is an incomplete proxy for clinical utility. The cosine learning rate schedule with 5\% warmup
produced smooth convergence without instability.

\begin{figure}[htbp]
\centering
\includegraphics[width=\columnwidth]{figures/fig1_training_curves.png}
\caption{Training dynamics over 4{,}050 steps (3 epochs).
Top-left: training and validation loss. Top-right: perplexity.
Bottom-left: cosine learning rate schedule. Bottom-right:
gradient norm with smoothed overlay.}
\label{fig:training}
\end{figure}

\begin{figure}[htbp]
\centering
\includegraphics[width=\columnwidth]{figures/fig2_epoch_loss.png}
\caption{Per-epoch training loss and validation loss.
Epoch 1: 0.3498 / 0.369. Epoch 2: 0.3224 / 0.369.
Epoch 3: 0.5197.}
\label{fig:epoch}
\end{figure}

\subsection{Clinical Accuracy}

Table~\ref{tab:main_results} summarises the primary evaluation
results. Aletheia achieved a Top-1 diagnostic accuracy of 80.0\%
and Top-3 accuracy of 100.0\%, indicating that the correct
diagnosis appears within the model's three highest-ranked
suggestions in all ten clinical presentations — a perfect Top-3 score.

\begin{table}[h!]
\centering
\caption{Main Evaluation Results}
\label{tab:main_results}
\begin{tabular}{lcc}
\toprule
\textbf{Metric} & \textbf{Score} & \textbf{Category} \\
\midrule
\multicolumn{3}{l}{\textit{Clinical Accuracy}} \\
Top-1 Diagnosis Accuracy  & 0.8000 & A \\
Top-3 Diagnosis Accuracy  & 1.0000 & A \\
\midrule
\multicolumn{3}{l}{\textit{Language Quality}} \\
ROUGE-1 (F1)              & 0.3830 & B \\
ROUGE-2 (F1)              & --- & B \\
ROUGE-L (F1)              & --- & B \\
BERTScore-F1              & 0.9086 & B \\
METEOR                    & 0.4665 & B \\
\midrule
\multicolumn{3}{l}{\textit{Calibration}} \\
Expected Calibration Error (ECE) & 0.2750 & C \\
\bottomrule
\end{tabular}
\end{table}

\begin{figure}[htbp]
\centering
\includegraphics[width=0.8\columnwidth]{figures/fig02_topk_accuracy.png}
\caption{Top-1 (80.0\%) and Top-3 (100.0\%) diagnostic accuracy
across 10 clinical case categories.}
\label{fig:topk}
\end{figure}

\begin{figure}[htbp]
\centering
\includegraphics[width=\columnwidth]{figures/fig01_per_condition_f1.png}
\caption{Per-condition precision, F1, and recall across all
10 evaluation case categories.}
\label{fig:f1}
\end{figure}

\begin{figure}[htbp]
\centering
\includegraphics[width=\columnwidth]{figures/fig03_confusion_matrix.png}
\caption{Confusion matrix for Top-1 differential diagnosis
predictions across 10 clinical conditions.}
\label{fig:cm}
\end{figure}

\subsection{Severity-Stratified Performance}

Fig.~\ref{fig:severity} shows diagnostic accuracy stratified by
clinical severity. The model demonstrates its highest accuracy on
conditions classified as High and Moderate severity, reflecting
the higher representation of these conditions in the training
dataset. Performance on Critical conditions, while lower in
absolute Top-1 accuracy, benefits from the high Top-3 accuracy
of 100.0\%, meaning the correct diagnosis is always present
in the ranked output for clinician review.

\begin{figure}[htbp]
\centering
\includegraphics[width=0.8\columnwidth]{figures/fig04_severity_accuracy.png}
\caption{Top-1 diagnostic accuracy stratified by clinical
severity (Critical, High, Moderate, Low).}
\label{fig:severity}
\end{figure}

\subsection{Reasoning Task Performance}

Fig.~\ref{fig:reasoning} shows Top-1 accuracy across the five
evaluated reasoning task types. The model performs most strongly
on test recommendation and evidence update tasks, reflecting the
structured nature of these outputs. Initial differential diagnosis
and follow-up question generation show lower but clinically
meaningful accuracy.

\begin{figure}[htbp]
\centering
\includegraphics[width=\columnwidth]{figures/fig05_reasoning_type.png}
\caption{Top-1 accuracy across five clinical reasoning task types.}
\label{fig:reasoning}
\end{figure}

\subsection{Language Quality}

Fig.~\ref{fig:language} shows the language quality metrics.
BERTScore-F1 of 0.909 indicates high semantic similarity between
model outputs and reference clinical answers. The relatively lower
ROUGE-1 score of 0.383 is expected for generative clinical
reasoning tasks where exact lexical overlap is a less meaningful
measure than semantic fidelity.

\begin{figure}[htbp]
\centering
\includegraphics[width=\columnwidth]{figures/fig06_language_quality.png}
\caption{Language quality metrics: ROUGE-1, ROUGE-2, ROUGE-L,
BERTScore-F1, and METEOR.}
\label{fig:language}
\end{figure}

\subsection{Calibration}

Fig.~\ref{fig:calibration} shows the reliability diagram and
Brier scores by severity. The ECE of 0.275 indicates moderate
calibration --- the model tends to assign probability estimates
that are directionally correct but somewhat overconfident on
conditions where it performs well. Calibration is worst for
Critical-severity conditions, which is the expected pattern for
models trained predominantly on synthetic data without
clinician-verified probability distributions.

\begin{figure}[htbp]
\centering
\includegraphics[width=\columnwidth]{figures/fig07_calibration.png}
\caption{Left: Reliability diagram with ECE~=~0.275.
Right: Brier score stratified by clinical severity.}
\label{fig:calibration}
\end{figure}

\subsection{Baseline Comparison}

Fig.~\ref{fig:baseline} compares fine-tuned Aletheia against the
unmodified Qwen2.5-3B-Instruct base model (zero-shot). Fine-tuning
produces a substantial improvement in Top-1 diagnostic accuracy
and ROUGE-1, confirming that the clinical reasoning capabilities
of Aletheia are a direct product of the fine-tuning process rather
than pre-existing in the base model.

\begin{figure}[htbp]
\centering
\includegraphics[width=\columnwidth]{figures/fig08_baseline_comparison.png}
\caption{Comparison of fine-tuned Aletheia vs base
Qwen2.5-3B-Instruct (zero-shot) on Top-1 accuracy and ROUGE-1.}
\label{fig:baseline}
\end{figure}

\begin{figure}[htbp]
\centering
\includegraphics[width=\columnwidth]{figures/fig09_rouge_distribution.png}
\caption{Box plot distributions of ROUGE-1, ROUGE-2, and ROUGE-L
scores across all evaluation cases, illustrating score spread
and consistency.}
\label{fig:rouge_dist}
\end{figure}

\subsection{Hardware Compliance}

Table~\ref{tab:efficiency} reports the deployment efficiency
metrics. Both GGUF quantisation formats pass the ADTC 2026 memory
ceiling of 7{,}168\,MB with substantial margin.

\begin{table}[h!]
\centering
\caption{Deployment Efficiency (ADTC 2026 Compliance)}
\label{tab:efficiency}
\begin{tabular}{lcc}
\toprule
\textbf{Metric} & \textbf{Q4\_K\_M} & \textbf{Q2\_K} \\
\midrule
Model file size           & 1.80\,GB & 1.19\,GB \\
Estimated peak RAM        & $\sim$3{,}630\,MB & $\sim$2{,}990\,MB \\
ADTC ceiling (7{,}168\,MB)& \textbf{PASS} & \textbf{PASS} \\
Margin below ceiling      & 3{,}538\,MB & 4{,}178\,MB \\
Internet required         & None & None \\
GPU required              & None & None \\
Context window            & 1{,}024 tokens & 1{,}024 tokens \\
\bottomrule
\end{tabular}
\end{table}

% ══════════════════════════════════════════════════════════════════
\section{Discussion}
\label{sec:discussion}

\subsection{Clinical Relevance of Results}

A Top-3 accuracy of 100.0\% is the most clinically meaningful
result from this evaluation. In practice, Aletheia presents the
clinician with a ranked list of differential diagnoses rather than
a single answer. A clinical officer reviewing a ranked list of
three diagnoses and finding the correct one present in 90\% of
cases represents a genuine decision-support capability --- the
system surfaces the diagnostic space that the clinician should be
considering, reducing the cognitive load of generating a
differential from scratch under time pressure.

The BERTScore-F1 of 0.909 indicates that the model's clinical
reasoning text is semantically faithful to expert reference answers.
This is important beyond accuracy scores: a system that names the
correct diagnosis but provides incorrect reasoning is clinically
dangerous. Aletheia's high BERTScore suggests that its explanations
are not only correct in conclusion but coherent in reasoning.

\subsection{Calibration and Clinical Safety}

The ECE of 0.275 indicates that the model's probability estimates
require careful interpretation. In clinical practice, the
probability values output by Aletheia should be treated as
\textit{relative rankings} rather than absolute probability
estimates. We recommend that the clinical interface present these
as ``more likely'' / ``less likely'' qualitative rankings rather
than specific percentages, until calibration can be improved
through clinician-validated feedback data.

\subsection{Comparison to Related Work}

Direct comparison to prior work is constrained by the lack of
standardised African clinical benchmarks. Med-PaLM\,2
\cite{singhal2023medpalm2} achieved 86.5\% on USMLE questions but
requires cloud inference and has no offline deployment pathway.
On the MedQA benchmark, Aletheia's base model (Qwen2.5-3B)
achieves competitive performance for its parameter count, and
fine-tuning produces substantial gains on Africa-specific
clinical scenarios that are not represented in standard benchmarks.

The key differentiator of Aletheia is not raw accuracy but
\textit{deployability}: the ability to run on a 8\,GB laptop
with no internet, consuming under 4\,GB RAM, at a model size of
under 2\,GB. No prior medical LLM system has demonstrated this
combination of capabilities in resource-constrained African healthcare contexts
across the continent --- a gap relevant across the entire continent.

\subsection{Limitations}

Several limitations of this work should be acknowledged:

\begin{enumerate}
  \item \textbf{Synthetic training data}: The majority of training
        samples are synthetically generated from hand-crafted case
        templates. While carefully designed to reflect African
        clinical presentation patterns, the probability distributions
        assigned to differentials reflect the authors' clinical
        knowledge synthesis rather than empirical epidemiological
        data derived from prospective patient cohorts. Ongoing
        clinician involvement from co-authors P.~B.~Kasasira and
        C.~B.~Okoboi is actively informing refinements to the
        case templates and probability estimates in preparation
        for the next training iteration.

  \item \textbf{Evaluation case set size}: The quantitative
        evaluation reported in this paper was conducted on 10 core
        case categories using automated metrics. A broader and
        clinician-graded evaluation set is currently under
        development as part of the ongoing validation activity
        with the School of Health Sciences co-authors, and will
        be reported in a subsequent study.

  \item \textbf{Calibration}: An ECE of 0.275 indicates that the
        model's probability estimates require improvement before
        they can be presented as clinically actionable confidence
        values. Initial feedback from clinician co-authors suggests
        that the directional ranking of differentials is clinically
        plausible in the majority of cases evaluated, though
        exact probability estimates should be treated as relative
        rankings rather than absolute likelihoods. Calibration
        improvement is a priority for the next training cycle.

  \item \textbf{Clinical validation scope}: An initial clinical
        evaluation of Aletheia is currently underway with two
        co-authors from the School of Health Sciences, Soroti
        University --- P.~B.~Kasasira and C.~B.~Okoboi --- who
        are practising clinicians actively testing and evaluating
        the system against real clinical presentations encountered
        in their practice. Early findings are encouraging: the
        ranked differential output has been rated as clinically
        plausible and useful for decision support in the majority
        of cases reviewed to date. However, this initial evaluation
        is limited in scale and has not yet been conducted under
        a formal study protocol with a defined patient cohort.
        A larger multi-site validation study is planned as the
        immediate next phase, involving clinical officers and
        physicians across district hospitals, regional referral
        hospitals, and primary health centres in Eastern Uganda,
        with a target enrolment of 500+ patient cases. This study
        will assess diagnostic concordance, clinical usability,
        acceptance, and safety, and is subject to IRB approval
        from Soroti University and the Uganda National Council
        for Science and Technology (UNCST).
\end{enumerate}

\subsection{Future Work}

Building on the current state of development and the ongoing
clinical evaluation, the following directions are planned:

\begin{itemize}
  \item Completion and formal reporting of the ongoing clinical
        evaluation with P.~B.~Kasasira and C.~B.~Okoboi, including
        a structured assessment of diagnostic concordance, usability,
        and clinician acceptance across a defined case set;

  \item Expansion to a multi-site prospective validation study
        across district hospitals, regional referral hospitals, and
        primary health centres in Eastern Uganda, targeting 500+
        patient cases and 50+ clinical officers, subject to IRB
        approval from Soroti University and the Uganda National
        Council for Science and Technology (UNCST);

  \item Incorporation of clinician feedback from the ongoing
        evaluation into the next training iteration, with particular
        focus on improving calibration (ECE), refining probability
        estimates, and expanding condition coverage from 50 to
        100+ clinical conditions;

  \item Development of a lightweight desktop application to
        complement the existing Gradio web UI and terminal CLI,
        making Aletheia accessible to non-technical clinical users;

  \item Kiswahili and Ateso language support to serve clinical
        officers across Eastern Uganda and the wider East African
        region (Section~\ref{sec:kiswahili_dataset});

  \item Submission to the Uganda National Drug Authority (NDA)
        under the software-as-a-medical-device regulatory pathway;

  \item Commercialisation through Arapai Technologies International Limited, with a distribution model targeting district
        health facilities across Uganda and the wider East African
        region at sustainable cost.
\end{itemize}

\subsection{Kiswahili Dataset Construction Strategy}
\label{sec:kiswahili_dataset}

Extending Aletheia to Kiswahili requires training data in a
symptoms-to-differential-diagnosis structured format consistent
with the schema used throughout this work. A survey of existing
public resources informed the decision to construct a dedicated
synthetic Kiswahili dataset rather than adopt or adapt an existing
corpus, for four reasons.

First, no openly available Kiswahili clinical dataset matches the
structured reasoning format required by Aletheia. General-purpose
Kiswahili corpora exist at scale --- including a corpus of
1.69~million sentences spanning a health (\textit{afya}) category
\cite{swahili_corpus_2024} --- but these consist of unstructured
text scraped from public websites rather than paired
symptom-to-diagnosis reasoning chains, and would require complete
reformatting before any clinical utility could be extracted.

Second, the closest comparable efforts in the literature are either
not yet publicly released or address a different problem framing.
A retrieval-augmented benchmarking methodology grounded in Kenyan
national clinical guidelines has been proposed to generate bilingual
English-Kiswahili clinical question-answer pairs validated by
practising Kenyan physicians \cite{alama_health_qa}, but produces
multiple-choice benchmark items rather than the ranked differential
diagnosis, investigation recommendation, and management-planning
schema that Aletheia requires. Similarly, a maternal health helpdesk
dataset of over 107{,}000 clinically labelled questions in Kiswahili
and Sheng has been deployed in Kenya \cite{jacaranda_health_2023},
but is both proprietary and scoped specifically to maternal health
triage rather than the broader 50-condition reasoning task addressed
by Aletheia.

Third, a recent scoping review of natural language processing for
public health in Africa confirms that language coverage in this
space remains uneven, with Kiswahili NLP resources for clinical
applications still largely in prototyping phases across the
literature \cite{nlp_africa_scoping_2025}. This indicates that the
absence of a ready-made dataset reflects a genuine gap in the field
rather than an oversight in our search.

Fourth, constructing a purpose-built synthetic dataset allows direct
alignment with Aletheia's existing 50-condition taxonomy, severity
classifications, and eight reasoning task types, ensuring
consistency with the English-language training data already
validated in Sections~\ref{sec:dataset}--\ref{sec:results}. This
approach mirrors the synthetic dataset construction strategy used
for the English-language component of this work, and is being
developed with direct input from clinician co-authors
P.~B.~Kasasira and C.~B.~Okoboi to ensure the Kiswahili symptom
vocabulary and clinical reasoning patterns remain faithful to
local practice.

An initial Kiswahili dataset covering all 50 conditions across
the eight reasoning task types is under construction, following
the same case-template methodology described in
Section~\ref{sec:dataset}, and will be incorporated into a
continued fine-tuning run from the existing LoRA adapter prior to
the prospective clinical validation study.

% ══════════════════════════════════════════════════════════════════
\section{Conclusion}
\label{sec:conclusion}

This paper has presented Aletheia, an offline-first clinical
decision support system designed for frontline healthcare workers
in sub-Saharan Africa. By combining QLoRA fine-tuning of a 3B
parameter language model on an Africa-weighted clinical reasoning dataset
(with particular emphasis on East African disease burden) with GGUF quantisation for edge deployment,
Aletheia achieves a Top-3 diagnostic accuracy of 100.0\% and
BERTScore-F1 of 0.909 while operating entirely within the memory
constraints of a standard 8\,GB laptop without internet
connectivity.

The results demonstrate that clinically meaningful AI-assisted
diagnostic reasoning is achievable at the primary care level in
resource-constrained settings without cloud infrastructure. In
the context of Uganda's physician-to-patient ratio of 1:25{,}000,
a system that presents the correct diagnosis in its top suggestion
in 80\% of cases, and within its top three suggestions in 100\% of cases, --- running on hardware already present in most
health facilities --- represents a meaningful force multiplier for
the clinical workforce.

The source code, inference scripts, and deployment tools for
Aletheia are publicly available at
\url{https://github.com/JosephWalusimbi-eng/Aletheia}.

Aletheia is released as open-source software under a public
health use licence \cite{aletheia_github}. The authors invite collaboration from Ugandan
clinicians, health informaticists, and the Ministry of Health for
the planned prospective validation study.

% ══════════════════════════════════════════════════════════════════
\section*{Acknowledgements}

The authors thank the clinical staff of Soroti Regional Referral
Hospital, Soroti District, whose daily practice under resource
constraints shed light on the problem this work addresses. The authors acknowledge the Africa Deep
Tech Challenge 2026 (ADTC 2026) for establishing the on-device
AI benchmark standard that shaped the hardware compliance
requirements of this work, and for providing a clear deployment
target that grounds this research in real-world constraints
relevant to the African technology ecosystem.

\textit{\textit{Conflict of Interest}: J.~Walusimbi and A.~M.~Oguti are
co-founders and directors of Arapai Technologies International Limited, a technology company based in Uganda. Aletheia is
intended for future commercialisation through this entity.
A.~Sserwadda, P.~B.~Kasasira, and C.~B.~Okoboi declare no
competing interests.}

% ══════════════════════════════════════════════════════════════════
\bibliographystyle{IEEEtran}
\begin{thebibliography}{99}
\sloppy
\raggedright

\bibitem{who2023workforce}
World Health Organization, ``Global Health Workforce Statistics,''
Geneva: WHO, 2023. [Online]. Available: \url{https://www.who.int/data/gho/data/themes/topics/health-workforce}

\bibitem{rajpurkar2022ai}
P. Rajpurkar, E. Chen, O. Banerjee, and E. J. Topol,
``AI in health and medicine,''
\textit{Nature Medicine}, vol. 28, no. 1, pp. 31--38, 2022.

\bibitem{esteva2017dermatologist}
A. Esteva \textit{et al.},
``Dermatologist-level classification of skin cancer with deep neural networks,''
\textit{Nature}, vol. 542, no. 7639, pp. 115--118, 2017.

\bibitem{mcdermott2021comprehensive}
M. B. A. McDermott \textit{et al.},
``A comprehensive EHR timeseries pre-training benchmark,''
in \textit{Proc. ACM Conf. Health, Inference, and Learning}, 2021, pp. 257--278.

\bibitem{nkosi2023challenges}
World Health Organization,
``Global Strategy on Digital Health 2020--2025,''
Geneva: WHO, 2021.
[Online]. Available: \url{https://www.who.int/docs/default-source/documents/gs4dhdaa2a9f352b0445bafbc79ca799dce4d.pdf}

\bibitem{dettmers2023qlora}
T. Dettmers, A. Pagnoni, A. Holtzman, and L. Zettlemoyer,
``QLoRA: Efficient finetuning of quantized LLMs,''
in \textit{Advances in Neural Information Processing Systems},
vol. 36, 2023.

\bibitem{gerganov2023llamacpp}
G. Gerganov, ``llama.cpp: Inference of LLaMA model in pure C/C++,''
GitHub, 2023. [Online]. Available: \url{https://github.com/ggerganov/llama.cpp}

\bibitem{openai2023gpt4}
OpenAI, ``GPT-4 Technical Report,'' arXiv:2303.08774, 2023.

\bibitem{singhal2023medpalm2}
K. Singhal \textit{et al.},
``Towards expert-level medical question answering with large language models,''
arXiv:2305.09617, 2023.

\bibitem{bolton2022biomedlm}
E. Bolton \textit{et al.},
``BioMedLM: A domain-specific large language model for biomedical text,''
Stanford CRFM, 2022.

\bibitem{huang2019clinicalbert}
K. Huang, J. Altosaar, and R. Ranganath,
``ClinicalBERT: Modeling clinical notes and predicting hospital readmission,''
arXiv:1904.05342, 2019.

\bibitem{fraser2007clinician}
H. S. F. Fraser, P. Biondich, D. Moodley, S. Choi, B. W. Mamlin, and P. Szolovits,
``Implementing electronic medical record systems in developing countries,''
\textit{Informatics in Primary Care}, vol. 13, no. 2, pp. 83--95, 2005.

\bibitem{labrique2013mhealth}
A. B. Labrique \textit{et al.},
``mHealth innovations as health system strengthening tools,''
\textit{Global Health: Science and Practice}, vol. 1, no. 2, pp. 160--171, 2013.

\bibitem{bastawrous2015ehealth}
A. Bastawrous and M. J. Armstrong,
``Mobile health use in low-and high-income countries,''
\textit{Journal of the Royal Society of Medicine}, vol. 106, no. 4,
pp. 130--142, 2013.

\bibitem{ada2020}
Ada Health GmbH, ``Ada: Your personal health guide,'' 2020.
[Online]. Available: \url{https://ada.com}

\bibitem{babylon2020}
Babylon Health, ``Babylon health AI,'' 2020.
[Online]. Available: \url{https://www.babylonhealth.com}

\bibitem{frantar2022gptq}
E. Frantar, S. Ashkboos, T. Hoefler, and D. Alistarh,
``GPTQ: Accurate post-training quantization for generative pre-trained transformers,''
arXiv:2210.17323, 2022.

\bibitem{lin2023awq}
J. Lin \textit{et al.},
``AWQ: Activation-aware weight quantization for LLM compression and acceleration,''
arXiv:2306.00978, 2023.

\bibitem{afrihealth2023}
World Health Organization Regional Office for Africa,
``Health Statistics and Information Systems: African Region,''
Brazzaville: WHO AFRO, 2023.
[Online]. Available: \url{https://www.afro.who.int/health-topics/health-statistics}

\bibitem{jin2021medqa}
D. Jin \textit{et al.},
``What disease does this patient have? A large-scale open domain question answering dataset from medical exams,''
\textit{Applied Sciences}, vol. 11, no. 14, p. 6421, 2021.

\bibitem{pmlr-v174-pal22a}
A. Pal \textit{et al.},
``MedMCQA: A large-scale multi-subject multi-choice dataset for medical domain question answering,''
in \textit{Proc. Conference on Health, Inference, and Learning},
PMLR, vol. 174, pp. 248--260, 2022.

\bibitem{qwen2025}
Qwen Team, ``Qwen2.5 Technical Report,'' arXiv:2412.15115, 2024.

\bibitem{peft2023}
S. Mangrulkar \textit{et al.},
``PEFT: State-of-the-art parameter-efficient fine-tuning methods,''
GitHub, 2022. [Online]. Available: \url{https://github.com/huggingface/peft}

\bibitem{lin2004rouge}
C.-Y. Lin,
``ROUGE: A package for automatic evaluation of summaries,''
in \textit{Text Summarization Branches Out}, 2004, pp. 74--81.

\bibitem{zhang2019bertscore}
T. Zhang \textit{et al.},
``BERTScore: Evaluating text generation with BERT,''
in \textit{Proc. ICLR}, 2020.

\bibitem{banerjee2005meteor}
S. Banerjee and A. Lavie,
``METEOR: An automatic metric for MT evaluation with improved correlation with human judgments,''
in \textit{Proc. ACL Workshop on Intrinsic and Extrinsic Evaluation Measures}, 2005.

\bibitem{naeini2015ece}
M. P. Naeini, G. F. Cooper, and M. Hauskrecht,
``Obtaining well calibrated probabilities using Bayesian binning,''
in \textit{Proc. AAAI}, 2015, pp. 2901--2907.

\bibitem{brier1950brier}
G. W. Brier,
``Verification of forecasts expressed in terms of probability,''
\textit{Monthly Weather Review}, vol. 78, no. 1, pp. 1--3, 1950.

\bibitem{adtc2026}
Africa Deep Tech Challenge, ``ADTC 2026: On-Device Language Model
Benchmark for African Hardware Contexts,'' 2026.
[Online]. Available: \url{https://africadeeptech.org/challenge-2026}

\bibitem{aletheia_github}
J. Walusimbi \textit{et al.},
``Aletheia: Offline-First Clinical Decision Support AI
for Sub-Saharan Africa,''
GitHub, 2026. [Online]. Available:
\url{https://github.com/JosephWalusimbi-eng/Aletheia}

\bibitem{swahili_corpus_2024}
E. F. Mwakapeje \textit{et al.},
``In the heart of Swahili: An exploration of data collection
methods and corpus curation for natural language processing,''
\textit{Data in Brief}, 2024.
[Online]. Available: \url{https://pmc.ncbi.nlm.nih.gov/articles/PMC11372376/}

\bibitem{alama_health_qa}
Alama Health Collaborative,
``Retrieval-Augmented Clinical Benchmarking for Contextual
Model Testing in Kenyan Primary Care: A Methodology Paper,''
arXiv:2507.14615, 2025.

\bibitem{jacaranda_health_2023}
Jacaranda Health,
``Healthcare queries in low-resource and code-mixed languages:
A maternal health helpdesk dataset in Swahili and Sheng,''
in \textit{Practical ML for Developing Countries Workshop, ICLR},
2023.

\bibitem{nlp_africa_scoping_2025}
S. Obadiah \textit{et al.},
``Natural Language Processing Technologies for Public Health
in Africa: Scoping Review,''
\textit{Journal of Medical Internet Research}, vol. 27, e68720, 2025.

\end{thebibliography}

\end{document}
